\section{Theory}
\label{theory}

% Goal: give the reader the conceptual background needed to understand
% the project. Biological concepts first, then the computational tools.
% Implementation specifics belong in §3.

% ── 1. eCLIP and RNA-binding proteins ─────────────────────────────────
%   • RBPs regulate RNA post-transcriptionally (splicing, stability,
%     localisation, translation). Dysregulation is linked to cancer and
%     neurological disease.
%   • eCLIP (Van Nostrand et al. 2016): UV crosslink → immunoprecipitation
%     → sequencing. Reverse transcriptase truncates at the crosslink site,
%     so the 5′ read end marks the contact position at ~nucleotide
%     resolution. Two IP replicates + SMInput (size-matched input) control.
%   • The three RBPs studied, their canonical RNA motifs, and the two
%     cell lines (K562 leukemia, HepG2 liver). Brief table or inline list.

% ── 2. Peak calling with PureCLIP ─────────────────────────────────────
%   • Conceptual role: converts a pile-up of aligned reads into a discrete
%     set of binding sites. Necessary step between raw sequencing and
%     biological interpretation.
%   • PureCLIP mechanism: HMM over read 5′-start density; states
%     "crosslinked" vs. "background"; optionally uses SMInput as a
%     covariate to subtract non-specific signal.
%   • Why parameters matter: bandwidth controls smoothing of the signal;
%     merge distance controls how nearby sites are grouped. Wrong values
%     produce either too few (over-filtered) or too many (noisy) sites,
%     and the right values depend on the dataset.

% ── 3. What makes a binding-site set "good" ───────────────────────────
%   • Reproducibility: real binding recurs across independent replicates;
%     noise does not. Chance-correction is needed because broad intervals
%     overlap by chance even for random site sets.
%   • Motif enrichment: genuine sites should be enriched for the RBP's
%     known RNA recognition sequence relative to shuffled background.
%     PWMs (position weight matrices) provide a quantitative scoring model.
%   • Benchmark recall: good calls should recover the curated set of
%     strong regions from the ENCODE reference pipeline.

% ── 4. LLMs as reasoning-based optimizers ─────────────────────────────
%   • Large language models (LLMs) can incorporate domain knowledge
%     (priors about biology, known failure modes) into their proposals,
%     going beyond pure black-box function optimization.
%   • In an iterative loop, an LLM can read structured feedback (score
%     deltas per iteration) and reason about which parameter to change
%     next — analogous to a human expert tuning settings by hand, but
%     automated and consistent.

% ── 5. Bayesian optimization (Optuna / TPE) ───────────────────────────
%   • Bayesian optimization builds a probabilistic surrogate model of the
%     objective surface from previous evaluations and uses it to choose
%     the next point to evaluate (exploration vs. exploitation trade-off).
%   • Tree-structured Parzen Estimator (TPE) models the distribution of
%     good and bad parameter sets separately and samples from the ratio.
%     No domain knowledge needed; purely data-driven.
